

\section*{Abstract}\label{sec:abstract}

This thesis designs, implements, and evaluates a PC-hosted, self-hosted proof-of-concept based on n8n for an IoT-style workflow-to-action pipeline. The scenario uses a temperature event to select a simulated fan-on or fan-off action. The objective is to examine whether workflow automation can be combined with a middleware boundary, structured action messages, validation design, and measurable evidence in a focused event-action system. The implementation uses Python middleware as the boundary between workflow logic and action execution, shared JSON contracts for sensor and action payloads, and a documented validation and safety design for allowed, malformed, and risky action outputs. The evaluation compares a deterministic n8n workflow based on a threshold rule with an agent-enhanced workflow that produces contract-shaped action output. Results were collected using CSV-based evaluation, with additional Raspberry Pi middleware validation where n8n remained on the PC while the middleware and action endpoint ran on the Raspberry Pi. The final successful local workflow runs produced the expected actions for the defined high- and low-temperature cases. The agent-enhanced workflow also produced contract-shaped fan-action output for the same cases, but with higher latency than the deterministic baseline. Raspberry Pi middleware validation showed that the PC-hosted n8n workflow could call the Pi-hosted middleware endpoint over the network and receive the expected simulated fan responses. The evaluation is limited to the defined cases and simulated fan actions, but it shows that the architecture can be implemented, measured, and discussed as a clear workflow-to-action proof-of-concept.

Keywords: Workflow automation, n8n, Internet of Things, middleware, JSON contracts, Raspberry Pi